\subsection{Symbolic Execution}\label{sec:se}

\begin{algorithm}[t] \footnotesize
  \caption{Symbolic Execution}
  \label{alg:symbexec}
  \begin{algorithmic}[1]
    \Require Program ($\pgm$), budget ($\budget$), query data ($\data$)
    \Ensure Test cases ($\testcases$), new query data ($\newdatas$)
    \Procedure{{\symb}}{$\pgm, \budget, \data$}
    \State $\langle \testcases, \newdatas \rangle \gets \langle \emptyset, \emptyset \rangle$ 
    \State $\states \gets \myset{\ainitstate}$ \Comment{$\ainitstate=(\initstmt, \initmemory, true, \emptyset)$} 
    \Repeat
        \State $\astate \gets \Choose(\states)$ \Comment{$\astate=(\stmt, \memory, \pc, \assignments)$}
        \If{$\stmt$ is a conditional statement with predicate $\bc$}
            % \State $\langle \truesat, \trueassignments, \falsesat, \falseassignments \rangle \gets \Solve(\pc, \bc)$
            \State $\truesat, \trueassignments \gets \Solve(\pc \land \bc, \data)$
            \State $\falsesat, \falseassignments \gets \Solve(\pc \land \neg \bc, \data)$
        	\If{$\truesat$}
		       $\states \gets \states \cup \myset{(\stmt_{1}, \nmemory, \pc \land \bc, \trueassignments)}$
		    \EndIf
        	\If{$\falsesat$}
		       $\states \gets \states \cup \myset{(\stmt_{2}, \nmemory, \pc \land \neg \bc, \falseassignments)}$
		    \EndIf
            \State $\newdatas \gets \newdatas \cup \myset{(\pc \land \bc, \trueassignments), (\pc \land \neg \bc, \falseassignments)}$
        \ElsIf {$\stmt$ reaches a program exit}
            \State $\testcases \gets \testcases \cup \assignments$   \Comment{$\astate = (\stmt, \_, \_, \assignments)$} 
        \Else
            \State $\states \gets \states \cup \myset{\newstate}$   \Comment{$\newstate = (\stmt', \nmemory, \pc, \assignments)$}
        \EndIf
    \Until{$\budget$ expires (or $\states = \emptyset$)}
    % \State $\branches \gets \Runtc(\pgm, \testcases)$
    \State \Return $\testcases, \newdatas$
    \EndProcedure
\end{algorithmic}
\end{algorithm}

Symbolic execution generates test cases to increase code coverage by iteratively updating so-called ``states''. We define a state as a tuple consisting of four elements: (1) the next statement to be executed ($\stmt$), (2) a symbolic memory ($\memory$) that maps each program variable to a symbolic variable (e.g., $x\to\alpha$), (3) a path condition ($\pc$), represented as the conjunction of symbolic branch predicates accumulated along the execution path (e.g., $\alpha \le 10$), and (4) a set of assignments ($\assignments$), where each assignment maps a symbolic variable to a concrete value satisfying $\pc$ (e.g., $\alpha = 5$). If $\pc$ is satisfiable, $\assignments$ contains a satisfying assignment returned by the SMT solver; otherwise, $\assignments = \emptyset$.

Algorithm~\ref{alg:symbexec} presents a standard symbolic execution procedure. It takes three inputs: the target program ($\pgm$), a testing time budget ($\budget$), and a query data set ($\data$). The role of $\data$ will be discussed later; conventional symbolic execution techniques~\cite{Cadar2008,du2026fense,yao2025empc,lee2025topseed} do not use $\data$ and can therefore be viewed as operating with $\data=\emptyset$. At line~2, the algorithm initializes two empty sets: the set of generated test cases ($\testcases$) and the set of query data ($\data$). Each element in $\data$ is a pair consisting of a path condition and its satisfying assignments (e.g., ($\alpha \le 10$,\{$\alpha = 5$\})). At line~3, the set of states ($\states$) is initialized as a singleton containing the state ($\ainitstate$, $\initmemory$, $true$, $\emptyset$), which corresponds to the target program's entry state.

After initialization, the algorithm repeatedly selects a state and updates the set $\states$ until the testing budget is exhausted, while accumulating generated test cases ($\testcases$) and new query data ($\newdatas$) (lines~4--16). At line~5, the $\Choose$ function selects a state $\astate$ from $\states$, after which $\astate$ is removed from $\states$ (i.e., $\states \gets \states \setminus \myset{\astate}$). If the next statement of $\astate=(\stmt,\memory,\pc,\assignments)$ corresponds to a conditional statement with predicate $\bc$, the $\Solve$ function evaluates the satisfiability of the extended path conditions $\pc \wedge \bc$ and $\pc \wedge \neg\bc$ (lines~7--8). For each satisfiable path condition, Algorithm~\ref{alg:symbexec} creates the corresponding successor state and adds it to $\states$ at lines~9--10: $(\stmt_{\it 1}, \nmemory, \pc \wedge \bc, \trueassignments)$ or $(\stmt_{\it 2}, \nmemory, \pc \wedge \neg\bc, \falseassignments)$. The corresponding path conditions and their satisfying assignments are also added to $\newdatas$ (line~11). If the selected state reaches a program exit, the algorithm generates a test case using the assignment stored in $\assignments$ (line~13). When the testing budget is exhausted, the algorithm returns the generated test cases $\testcases$ and new query data $\newdatas$ (line~17).

\iffalse
Recent advancements in symbolic execution have significantly improved its effectiveness based on Algorithm~\ref{alg:symbexec}. For instance, search strategies~\cite{learch,yao2025empc}, corresponding to the $\Choose$ function at line~5, prioritize selecting the most promising states first from the set $\states$ based on predefined criteria. Conversely, state-pruning strategies~\cite{homi,chen2018dynamic} aim to eliminate redundant states from $\states$ by terminating them early. Some techniques---such as the $\Solve$ function at line~10---focus on efficiently reducing constraint solving costs~\cite{kleeaaqc,partialcaching}. Additionally, an orthogonal technique~\cite{symtuner} attempts to automatically tune external parameters of the symbolic execution engine (the $\symb$ procedure). These state-of-the-art techniques have significantly improved code coverage and facilitated the discovery of subtle bugs, compared to the generic symbolic execution algorithm (Algorithm~\ref{alg:symbexec}).

\myparagraph{\bf Modern symbolic execution}
Recent advances in symbolic execution have focused on optimizing each function in Algorithm~\ref{alg:symbexec} under a single budget. For example, heuristics~\cite{yao2025empc,featmaker} assign priorities to each state in the set $\states$ at line~5, guiding which state the $\Choose$ function explores first. State-pruning approaches~\cite{du2026fense,homi} filter out inefficient states from $\states$, allowing the Select function to explore promising states more quickly. At line~7, constraint-solving~\cite{symloc,mikek2024smt} techniques refine path condition $\pc$ to reduce the cost of computing satisfiability in the $\Solve$ function. Lastly, caching methods~\cite{shuai2024partial,klee_aaqc} store query results in $\data$ and reuse them to avoid unnecessary solver calls in the $\Solve$ function when the same constraints reach. All these techniques outperformed standard symbolic execution in terms of code coverage and bug finding.
\fi